% --- Archon physics formalization source begin ---
% archon:physics
% archon:covers PhyXMiniProblems/problem_phyx_mini_0316.lean
% archon:source-report reports/phyx_mini/problem_phyx_mini_0316.source.json

\chapter{Physics problem phyx\_mini\_0316}
\label{ch:PhyXMiniProblems_problem_phyx_mini_0316}

\paragraph{Problem source.}
\#\# Physical scenario

Figure shows two point sources \$S\_1\$ and \$S\_2\$ that emit sound of wavelength \$\textbackslash{}lambda = 2.00\$ m. The emissions are isotropic and in phase, and the separation between the sources is \$d = 16.0\$ m. At any point \$P\$ on the \$x\$ axis, the wave from \$S\_1\$ and the wave from \$S\_2\$ interfere.

\#\# Question

At what distance \$x\$ do the waves have a phase difference of \$1.50\textbackslash{}lambda\$?

\#\# Answer choices

- A: A: 38.6 m

- B: B: 39.8 m

- C: C: 40.6 m

- D: D: 41.2 m

\#\# Figure caption (auxiliary; use the image as primary evidence)

The image depicts a two-dimensional Cartesian coordinate system with the x-axis and y-axis. 

- There are two points marked on the y-axis: \textbackslash{}( S\_1 \textbackslash{}) and \textbackslash{}( S\_2 \textbackslash{}). 
  - \textbackslash{}( S\_1 \textbackslash{}) is above \textbackslash{}( S\_2 \textbackslash{}), and there's a vertical distance \textbackslash{}( d \textbackslash{}) between them.
  - Both points \textbackslash{}( S\_1 \textbackslash{}) and \textbackslash{}( S\_2 \textbackslash{}) are marked with red dots.

- A point \textbackslash{}( P \textbackslash{}) is located on the positive x-axis, marked with a black dot.

- The axes intersect at the origin, and the axes are labeled with \textbackslash{}( x \textbackslash{}) and \textbackslash{}( y \textbackslash{}).

- The text labels include \textbackslash{}( S\_1 \textbackslash{}), \textbackslash{}( S\_2 \textbackslash{}), \textbackslash{}( P \textbackslash{}), \textbackslash{}( x \textbackslash{}), \textbackslash{}( y \textbackslash{}), and \textbackslash{}( d \textbackslash{}).

\#\# Dataset metadata

Resonance and Harmonics; Spatial Relation Reasoning, Multi-Formula Reasoning

\paragraph{Recorded answer/context.}
D: D: 41.2 m

\paragraph{Figure/image path.}
/root/proposal\_for\_physic/science-mango/phyx\_mini\_run/phyx\_data/test\_image/316.png

\paragraph{Formalization target.}
create a compiling Lean file with sorry bodies at `PhyXMiniProblems/problem\_phyx\_mini\_0316.lean`. The Lean declarations must preserve the physical quantities, dimensions or dimensional roles, figure labels, governing-law hypotheses, and final relation expressed by this problem.
Use Mathlib/Physlib names found through LeanExplore where available. If a physics API is missing, introduce faithful local abstractions rather than scalar placeholder aliases.

\begin{theorem}[Physics formalization target]
\label{thm:physics:phyx_mini_0316:target}
\lean{PhyXMiniProblems.ProblemPhyXMini0316.problem_phyx_mini_0316}
\uses{def:physics:phyx-mini-0316:phyxminiproblems-problemphyxmini0316-twopointsourceinterferencesetup, def:physics:phyx-mini-0316:phyxminiproblems-problemphyxmini0316-matchescoherentinphasesourcedescription, def:physics:phyx-mini-0316:phyxminiproblems-problemphyxmini0316-matchesproblemandquestionreadouts, def:physics:phyx-mini-0316:phyxminiproblems-problemphyxmini0316-matchessuppliedfigure, def:physics:phyx-mini-0316:phyxminiproblems-problemphyxmini0316-hasphysicalacousticparameters, def:physics:phyx-mini-0316:phyxminiproblems-problemphyxmini0316-satisfiesperpendicularpathgeometry, def:physics:phyx-mini-0316:phyxminiproblems-problemphyxmini0316-satisfiesgeometricpathdifferencelaw, def:physics:phyx-mini-0316:phyxminiproblems-problemphyxmini0316-satisfiesinphaseacousticphaselaw, def:physics:phyx-mini-0316:phyxminiproblems-problemphyxmini0316-requestedphasedistancesetmeters, def:physics:phyx-mini-0316:phyxminiproblems-problemphyxmini0316-isuniquematchinganswer}
The assigned autoformalize agent should translate this physics problem into Lean declarations in the covered file, with theorem and lemma proof bodies written as `by sorry`.
\end{theorem}
\begin{proof}
This is an autoformalization task, not a proof task. Produce statements that can later be proved by the physics prover without weakening the physical model.
\end{proof}

% --- Archon declaration topology begin ---
\section{Declaration topology}
\begin{definition}[Length Quantity]
  \label{def:physics:phyx-mini-0316:phyxminiproblems-problemphyxmini0316-lengthquantity}
  \lean{PhyXMiniProblems.ProblemPhyXMini0316.LengthQuantity}
  A nonnegative physical length, independent of the unit used to read it
\end{definition}
\begin{proof}
  This is the declared model definition of Length Quantity.
\end{proof}
\begin{definition}[Signed Length Quantity]
  \label{def:physics:phyx-mini-0316:phyxminiproblems-problemphyxmini0316-signedlengthquantity}
  \lean{PhyXMiniProblems.ProblemPhyXMini0316.SignedLengthQuantity}
  A signed physical length used for Cartesian coordinates in the figure
\end{definition}
\begin{proof}
  This is the declared model definition of Signed Length Quantity.
\end{proof}
\begin{definition}[length Readout]
  \label{def:physics:phyx-mini-0316:phyxminiproblems-problemphyxmini0316-lengthreadout}
  \lean{PhyXMiniProblems.ProblemPhyXMini0316.lengthReadout}
  \uses{def:physics:phyx-mini-0316:phyxminiproblems-problemphyxmini0316-lengthquantity}
  Read a nonnegative physical length in a selected length unit
\end{definition}
\begin{proof}
  This is the declared model definition of length Readout.
\end{proof}
\begin{definition}[signed Length Readout]
  \label{def:physics:phyx-mini-0316:phyxminiproblems-problemphyxmini0316-signedlengthreadout}
  \lean{PhyXMiniProblems.ProblemPhyXMini0316.signedLengthReadout}
  \uses{def:physics:phyx-mini-0316:phyxminiproblems-problemphyxmini0316-signedlengthquantity}
  Read a signed Cartesian coordinate in a selected length unit
\end{definition}
\begin{proof}
  This is the declared model definition of signed Length Readout.
\end{proof}
\begin{definition}[length In Meters]
  \label{def:physics:phyx-mini-0316:phyxminiproblems-problemphyxmini0316-lengthinmeters}
  \lean{PhyXMiniProblems.ProblemPhyXMini0316.lengthInMeters}
  \uses{def:physics:phyx-mini-0316:phyxminiproblems-problemphyxmini0316-lengthquantity, def:physics:phyx-mini-0316:phyxminiproblems-problemphyxmini0316-lengthreadout}
  Meter readout of a nonnegative physical length
\end{definition}
\begin{proof}
  This is the declared model definition of length In Meters.
\end{proof}
\begin{definition}[signed Length In Meters]
  \label{def:physics:phyx-mini-0316:phyxminiproblems-problemphyxmini0316-signedlengthinmeters}
  \lean{PhyXMiniProblems.ProblemPhyXMini0316.signedLengthInMeters}
  \uses{def:physics:phyx-mini-0316:phyxminiproblems-problemphyxmini0316-signedlengthquantity, def:physics:phyx-mini-0316:phyxminiproblems-problemphyxmini0316-signedlengthreadout}
  Meter readout of a signed figure coordinate
\end{definition}
\begin{proof}
  This is the declared model definition of signed Length In Meters.
\end{proof}
\begin{definition}[Source Label]
  \label{def:physics:phyx-mini-0316:phyxminiproblems-problemphyxmini0316-sourcelabel}
  \lean{PhyXMiniProblems.ProblemPhyXMini0316.SourceLabel}
  Labels of the two sound sources printed in the primary figure
\end{definition}
\begin{proof}
  This is the declared model definition of Source Label.
\end{proof}
\begin{definition}[Figure Point Label]
  \label{def:physics:phyx-mini-0316:phyxminiproblems-problemphyxmini0316-figurepointlabel}
  \lean{PhyXMiniProblems.ProblemPhyXMini0316.FigurePointLabel}
  Point labels printed in the primary figure
\end{definition}
\begin{proof}
  This is the declared model definition of Figure Point Label.
\end{proof}
\begin{definition}[Coordinate Axis]
  \label{def:physics:phyx-mini-0316:phyxminiproblems-problemphyxmini0316-coordinateaxis}
  \lean{PhyXMiniProblems.ProblemPhyXMini0316.CoordinateAxis}
  Cartesian axis labels printed in the primary figure
\end{definition}
\begin{proof}
  This is the declared model definition of Coordinate Axis.
\end{proof}
\begin{definition}[Acoustic Source Kind]
  \label{def:physics:phyx-mini-0316:phyxminiproblems-problemphyxmini0316-acousticsourcekind}
  \lean{PhyXMiniProblems.ProblemPhyXMini0316.AcousticSourceKind}
  The source type stated in the problem
\end{definition}
\begin{proof}
  This is the declared model definition of Acoustic Source Kind.
\end{proof}
\begin{definition}[Acoustic Medium]
  \label{def:physics:phyx-mini-0316:phyxminiproblems-problemphyxmini0316-acousticmedium}
  \lean{PhyXMiniProblems.ProblemPhyXMini0316.AcousticMedium}
  The common homogeneous propagation medium implicit in the sound diagram
\end{definition}
\begin{proof}
  This is the declared model definition of Acoustic Medium.
\end{proof}
\begin{definition}[Plane Point]
  \label{def:physics:phyx-mini-0316:phyxminiproblems-problemphyxmini0316-planepoint}
  \lean{PhyXMiniProblems.ProblemPhyXMini0316.PlanePoint}
  \uses{def:physics:phyx-mini-0316:phyxminiproblems-problemphyxmini0316-signedlengthquantity}
  A Cartesian point whose coordinates are dimensionful signed lengths
\end{definition}
\begin{proof}
  This is the declared model definition of Plane Point.
\end{proof}
\begin{definition}[coordinate In Meters]
  \label{def:physics:phyx-mini-0316:phyxminiproblems-problemphyxmini0316-coordinateinmeters}
  \lean{PhyXMiniProblems.ProblemPhyXMini0316.coordinateInMeters}
  \uses{def:physics:phyx-mini-0316:phyxminiproblems-problemphyxmini0316-signedlengthinmeters, def:physics:phyx-mini-0316:phyxminiproblems-problemphyxmini0316-coordinateaxis, def:physics:phyx-mini-0316:phyxminiproblems-problemphyxmini0316-planepoint}
  Meter readout of one coordinate of a physical figure point
\end{definition}
\begin{proof}
  This is the declared model definition of coordinate In Meters.
\end{proof}
\begin{definition}[Two Point Source Interference Setup]
  \label{def:physics:phyx-mini-0316:phyxminiproblems-problemphyxmini0316-twopointsourceinterferencesetup}
  \lean{PhyXMiniProblems.ProblemPhyXMini0316.TwoPointSourceInterferenceSetup}
  \uses{def:physics:phyx-mini-0316:phyxminiproblems-problemphyxmini0316-lengthquantity, def:physics:phyx-mini-0316:phyxminiproblems-problemphyxmini0316-sourcelabel, def:physics:phyx-mini-0316:phyxminiproblems-problemphyxmini0316-figurepointlabel, def:physics:phyx-mini-0316:phyxminiproblems-problemphyxmini0316-acousticsourcekind, def:physics:phyx-mini-0316:phyxminiproblems-problemphyxmini0316-acousticmedium, def:physics:phyx-mini-0316:phyxminiproblems-problemphyxmini0316-planepoint}
  The physical quantities, figure-labelled objects, paths, and phase observable. observationPointPAtMeters x is the point labelled P whose nonnegative horizontal coordinate is read as x meters. The path lengths and their geometric difference remain dimensionful. No field assigns a numerical value to the requested coordinate x or selects an answer choice
\end{definition}
\begin{proof}
  This is the declared model definition of Two Point Source Interference Setup.
\end{proof}
\begin{definition}[Matches Coherent In Phase Source Description]
  \label{def:physics:phyx-mini-0316:phyxminiproblems-problemphyxmini0316-matchescoherentinphasesourcedescription}
  \lean{PhyXMiniProblems.ProblemPhyXMini0316.MatchesCoherentInPhaseSourceDescription}
  \uses{def:physics:phyx-mini-0316:phyxminiproblems-problemphyxmini0316-sourcelabel, def:physics:phyx-mini-0316:phyxminiproblems-problemphyxmini0316-twopointsourceinterferencesetup}
  The qualitative acoustic description: the sources are isotropic, mutually coherent, and initially in phase in one common medium. Isotropic propagation also makes both waves available to overlap at every pictured point on the nonnegative x axis. This contains no distinguished distance
\end{definition}
\begin{proof}
  This is the declared model definition of Matches Coherent In Phase Source Description.
\end{proof}
\begin{definition}[Matches Problem And Question Readouts]
  \label{def:physics:phyx-mini-0316:phyxminiproblems-problemphyxmini0316-matchesproblemandquestionreadouts}
  \lean{PhyXMiniProblems.ProblemPhyXMini0316.MatchesProblemAndQuestionReadouts}
  \uses{def:physics:phyx-mini-0316:phyxminiproblems-problemphyxmini0316-lengthinmeters, def:physics:phyx-mini-0316:phyxminiproblems-problemphyxmini0316-twopointsourceinterferencesetup}
  Numerical data stated in the problem and question: lambda = 2.00 m, d = 16.0 m, and the requested phase difference is 1.50 wavelengths. No observation distance or answer-choice value occurs in these readouts
\end{definition}
\begin{proof}
  This is the declared model definition of Matches Problem And Question Readouts.
\end{proof}
\begin{definition}[Matches Supplied Figure]
  \label{def:physics:phyx-mini-0316:phyxminiproblems-problemphyxmini0316-matchessuppliedfigure}
  \lean{PhyXMiniProblems.ProblemPhyXMini0316.MatchesSuppliedFigure}
  \uses{def:physics:phyx-mini-0316:phyxminiproblems-problemphyxmini0316-lengthinmeters, def:physics:phyx-mini-0316:phyxminiproblems-problemphyxmini0316-coordinateinmeters, def:physics:phyx-mini-0316:phyxminiproblems-problemphyxmini0316-twopointsourceinterferencesetup}
  Coordinate and label information read from the primary image. S₁ is at the axes' intersection, S₂ is directly below it by d, and the point P runs along the nonnegative x axis. These are figure readouts only; they contain no phase condition
\end{definition}
\begin{proof}
  This is the declared model definition of Matches Supplied Figure.
\end{proof}
\begin{definition}[Has Physical Acoustic Parameters]
  \label{def:physics:phyx-mini-0316:phyxminiproblems-problemphyxmini0316-hasphysicalacousticparameters}
  \lean{PhyXMiniProblems.ProblemPhyXMini0316.HasPhysicalAcousticParameters}
  \uses{def:physics:phyx-mini-0316:phyxminiproblems-problemphyxmini0316-lengthinmeters, def:physics:phyx-mini-0316:phyxminiproblems-problemphyxmini0316-sourcelabel, def:physics:phyx-mini-0316:phyxminiproblems-problemphyxmini0316-twopointsourceinterferencesetup}
  Positivity and nondegeneracy of the physical wave and path quantities
\end{definition}
\begin{proof}
  This is the declared model definition of Has Physical Acoustic Parameters.
\end{proof}
\begin{definition}[Satisfies Perpendicular Path Geometry]
  \label{def:physics:phyx-mini-0316:phyxminiproblems-problemphyxmini0316-satisfiesperpendicularpathgeometry}
  \lean{PhyXMiniProblems.ProblemPhyXMini0316.SatisfiesPerpendicularPathGeometry}
  \uses{def:physics:phyx-mini-0316:phyxminiproblems-problemphyxmini0316-lengthinmeters, def:physics:phyx-mini-0316:phyxminiproblems-problemphyxmini0316-twopointsourceinterferencesetup}
  Straight-line Euclidean propagation geometry on the horizontal observation ray. A point x meters from S₁ is x meters from S₁ and sqrt (x\textasciicircum{}2 + d\textasciicircum{}2) meters from S₂. This law is quantified over arbitrary nonnegative x and does not single out the requested phase condition
\end{definition}
\begin{proof}
  This is the declared model definition of Satisfies Perpendicular Path Geometry.
\end{proof}
\begin{definition}[Satisfies Geometric Path Difference Law]
  \label{def:physics:phyx-mini-0316:phyxminiproblems-problemphyxmini0316-satisfiesgeometricpathdifferencelaw}
  \lean{PhyXMiniProblems.ProblemPhyXMini0316.SatisfiesGeometricPathDifferenceLaw}
  \uses{def:physics:phyx-mini-0316:phyxminiproblems-problemphyxmini0316-lengthreadout, def:physics:phyx-mini-0316:phyxminiproblems-problemphyxmini0316-twopointsourceinterferencesetup}
  In one homogeneous medium, the geometric propagation-path difference is the magnitude of the difference of the two source-to-observer path lengths. The relation is required in every length unit and contains no requested value
\end{definition}
\begin{proof}
  This is the declared model definition of Satisfies Geometric Path Difference Law.
\end{proof}
\begin{definition}[Satisfies In Phase Acoustic Phase Law]
  \label{def:physics:phyx-mini-0316:phyxminiproblems-problemphyxmini0316-satisfiesinphaseacousticphaselaw}
  \lean{PhyXMiniProblems.ProblemPhyXMini0316.SatisfiesInPhaseAcousticPhaseLaw}
  \uses{def:physics:phyx-mini-0316:phyxminiproblems-problemphyxmini0316-lengthreadout, def:physics:phyx-mini-0316:phyxminiproblems-problemphyxmini0316-twopointsourceinterferencesetup}
  For coherent in-phase sources, phase difference measured in wavelength cycles is geometric path difference divided by wavelength. Stating the ratio in every length unit records its dimensionless physical meaning. This is a general wave law and does not assert where a particular phase occurs
\end{definition}
\begin{proof}
  This is the declared model definition of Satisfies In Phase Acoustic Phase Law.
\end{proof}
\begin{definition}[Has Requested Phase Difference At Meters]
  \label{def:physics:phyx-mini-0316:phyxminiproblems-problemphyxmini0316-hasrequestedphasedifferenceatmeters}
  \lean{PhyXMiniProblems.ProblemPhyXMini0316.HasRequestedPhaseDifferenceAtMeters}
  \uses{def:physics:phyx-mini-0316:phyxminiproblems-problemphyxmini0316-twopointsourceinterferencesetup}
  A nonnegative meter coordinate at which the question's phase occurs
\end{definition}
\begin{proof}
  This is the declared model definition of Has Requested Phase Difference At Meters.
\end{proof}
\begin{definition}[requested Phase Distance Set Meters]
  \label{def:physics:phyx-mini-0316:phyxminiproblems-problemphyxmini0316-requestedphasedistancesetmeters}
  \lean{PhyXMiniProblems.ProblemPhyXMini0316.requestedPhaseDistanceSetMeters}
  \uses{def:physics:phyx-mini-0316:phyxminiproblems-problemphyxmini0316-twopointsourceinterferencesetup, def:physics:phyx-mini-0316:phyxminiproblems-problemphyxmini0316-hasrequestedphasedifferenceatmeters}
  The set of all nonnegative meter coordinates satisfying the question
\end{definition}
\begin{proof}
  This is the declared model definition of requested Phase Distance Set Meters.
\end{proof}
\begin{definition}[Answer Choice]
  \label{def:physics:phyx-mini-0316:phyxminiproblems-problemphyxmini0316-answerchoice}
  \lean{PhyXMiniProblems.ProblemPhyXMini0316.AnswerChoice}
  Labels of the four displayed multiple-choice answers
\end{definition}
\begin{proof}
  This is the declared model definition of Answer Choice.
\end{proof}
\begin{definition}[displayed Distance In Meters]
  \label{def:physics:phyx-mini-0316:phyxminiproblems-problemphyxmini0316-displayeddistanceinmeters}
  \lean{PhyXMiniProblems.ProblemPhyXMini0316.displayedDistanceInMeters}
  \uses{def:physics:phyx-mini-0316:phyxminiproblems-problemphyxmini0316-answerchoice}
  Distance in meters printed beside each answer label
\end{definition}
\begin{proof}
  This is the declared model definition of displayed Distance In Meters.
\end{proof}
\begin{definition}[Rounds To Displayed Distance]
  \label{def:physics:phyx-mini-0316:phyxminiproblems-problemphyxmini0316-roundstodisplayeddistance}
  \lean{PhyXMiniProblems.ProblemPhyXMini0316.RoundsToDisplayedDistance}
  \uses{def:physics:phyx-mini-0316:phyxminiproblems-problemphyxmini0316-answerchoice, def:physics:phyx-mini-0316:phyxminiproblems-problemphyxmini0316-displayeddistanceinmeters}
  An exact distance rounds to a displayed tenth-meter choice when it differs from that printed value by less than half of 0.1 m
\end{definition}
\begin{proof}
  This is the declared model definition of Rounds To Displayed Distance.
\end{proof}
\begin{definition}[Is Unique Matching Answer]
  \label{def:physics:phyx-mini-0316:phyxminiproblems-problemphyxmini0316-isuniquematchinganswer}
  \lean{PhyXMiniProblems.ProblemPhyXMini0316.IsUniqueMatchingAnswer}
  \uses{def:physics:phyx-mini-0316:phyxminiproblems-problemphyxmini0316-answerchoice, def:physics:phyx-mini-0316:phyxminiproblems-problemphyxmini0316-roundstodisplayeddistance}
  Exactly one displayed value agrees with the exact derived distance
\end{definition}
\begin{proof}
  This is the declared model definition of Is Unique Matching Answer.
\end{proof}
\begin{lemma}[requested Phase Distance Set eq singleton]
  \label{lem:physics:phyx-mini-0316:phyxminiproblems-problemphyxmini0316-requestedphasedistanceset-eq-singleton}
  \lean{PhyXMiniProblems.ProblemPhyXMini0316.requestedPhaseDistanceSet_eq_singleton}
  \uses{def:physics:phyx-mini-0316:phyxminiproblems-problemphyxmini0316-twopointsourceinterferencesetup, def:physics:phyx-mini-0316:phyxminiproblems-problemphyxmini0316-matchesproblemandquestionreadouts, def:physics:phyx-mini-0316:phyxminiproblems-problemphyxmini0316-hasphysicalacousticparameters, def:physics:phyx-mini-0316:phyxminiproblems-problemphyxmini0316-satisfiesperpendicularpathgeometry, def:physics:phyx-mini-0316:phyxminiproblems-problemphyxmini0316-satisfiesgeometricpathdifferencelaw, def:physics:phyx-mini-0316:phyxminiproblems-problemphyxmini0316-satisfiesinphaseacousticphaselaw, def:physics:phyx-mini-0316:phyxminiproblems-problemphyxmini0316-requestedphasedistancesetmeters}
  The perpendicular-path and in-phase wave laws determine the unique exact coordinate x = 247/6 m for a 3/2-wavelength phase difference. The exact value occurs only in this conclusion, not in setup data or governing laws
\end{lemma}
\begin{proof}
  The conclusion follows from the governing relations and figure readouts named in the statement of requested Phase Distance Set eq singleton.
\end{proof}
% --- Archon declaration topology end ---
% --- Archon physics formalization source end ---
